\documentclass[11pt]{article}

\usepackage[margin=1in]{geometry}
\usepackage{amsmath,amssymb}
\usepackage{array}
\usepackage{enumitem}

\title{Lecture 3 Notes: Plato's Rationalism, and Aristotle}
\author{Philosophy of Mathematics}
\date{}

\begin{document}

\maketitle

\section*{Central Theme}

Plato and Aristotle offer two rival ways to explain the exactness and objectivity of
mathematics.

\[
\boxed{
\text{Plato separates mathematics from the physical world.}
}
\]

\[
\boxed{
\text{Aristotle embeds mathematics in the physical world.}
}
\]

The guiding contrast:

\[
\begin{array}{c|c}
\textbf{Plato} & \textbf{Aristotle} \\
\hline
\text{explains exactness} & \text{explains applicability} \\
\text{makes application mysterious} & \text{makes exactness mysterious}
\end{array}
\]

\section*{1. Greek Mathematics: Exactness and Proof}

Greek mathematics introduced a distinctive demand for exactitude and rigorous proof.

The story of doubling the altar illustrates this demand. A practical problem becomes a
mathematical problem: not approximately double the altar, but exactly double it.

The resulting problem is the problem of doubling the cube:

\[
\text{Given a line segment } s,\text{ construct } x \text{ such that } x^3 = 2s^3.
\]

Similar ancient problems include:

\[
\text{trisecting an angle}
\]

\[
\text{squaring the circle}
\]

In each case, approximation was not enough. The Greek mathematicians wanted exact
construction and proof.

\subsection*{Teaching Point}

Greek mathematics is philosophically important because it treats exact proof, not practical
approximation, as the standard of mathematical success.

\subsection*{Whiteboard}

\[
\text{Practical problem}
\quad \longrightarrow \quad
\text{exact construction}
\quad \longrightarrow \quad
\text{proof}
\]

\[
2.0004 \neq 2
\quad
\text{for the mathematician}
\]

\section*{2. Plato's World of Being}

Plato begins from a gap between imperfect physical things and perfect intelligible ideals.

We encounter beautiful things, just actions, and circular objects. But we do not encounter
Beauty itself, Justice itself, or the perfect Circle in the physical world.

Plato explains this by distinguishing two realms:

\[
\begin{array}{c|c}
\textbf{World of Becoming} & \textbf{World of Being} \\
\hline
\text{physical} & \text{intelligible} \\
\text{changing} & \text{eternal} \\
\text{imperfect} & \text{perfect} \\
\text{known by senses} & \text{known by thought}
\end{array}
\]

The Forms belong to the world of Being. Physical objects resemble, participate in, or
approximate the Forms.

\subsection*{Teaching Point}

Plato's metaphysics explains why physical things are imperfect: they are not the standards
themselves, but only imperfect participants in those standards.

\subsection*{Whiteboard}

\[
\text{beautiful things} \neq \text{Beauty itself}
\]

\[
\text{just actions} \neq \text{Justice itself}
\]

\[
\text{drawn circles} \neq \text{Circle itself}
\]

\section*{3. Plato's Epistemology: Thought over Sense}

For Plato, the senses give us access to the world of Becoming. Reason gives us access to
the world of Being.

This is why mathematics is central. Mathematics trains the soul to turn away from the
changing world of sense and toward eternal, intelligible reality.

The \emph{Meno} provides an important example. Socrates guides a slave to a geometrical
truth without directly teaching it to him. Plato uses this as evidence for recollection: what
we call learning is really remembering.

Even if we do not accept the doctrine of recollection, the main point remains:
mathematical knowledge is not ordinary sensory knowledge.

\subsection*{Teaching Point}

For Plato, mathematics is a bridge from sensory experience to rational insight.

\subsection*{Whiteboard}

\[
\text{sense perception}
\quad \longrightarrow \quad
\text{Becoming}
\]

\[
\text{reason}
\quad \longrightarrow \quad
\text{Being}
\]

\[
\text{mathematics}
\quad =
\quad
\text{training for ascent}
\]

\section*{4. Geometry as Plato's Model}

Geometry gives Plato a clear example of the gap between the physical and the intelligible.

A theorem of geometry concerns perfect points, lines, and circles. But no physical drawing
contains a perfect point, a breadthless line, or a perfect circle.

Consider the theorem:

\[
\text{A tangent to a circle intersects the circle at exactly one point.}
\]

In a physical drawing, the line and the circle overlap in a small region. That does not
disconfirm the theorem. Instead, for Plato, it shows that the theorem is not about the
drawing.

\subsection*{Teaching Point}

A diagram is not the object of geometry. It is an image or aid that helps the mind grasp the
ideal geometrical object.

\subsection*{Whiteboard}

\[
\text{drawn circle}
\quad \neq \quad
\text{Circle itself}
\]

\[
\text{drawn tangent}
\quad \neq \quad
\text{ideal tangent}
\]

\[
\text{diagram} = \text{image, not object}
\]

\section*{5. Plato's Divided Line}

In the \emph{Republic}, Plato uses the divided line to represent levels of reality and cognition.

\[
\begin{array}{c}
\text{The Good} \\
\hline
\text{Forms} \\
\hline
\text{Mathematical objects} \\
\hline
\text{Physical objects} \\
\hline
\text{Reflections/images}
\end{array}
\]

Mathematical objects are above physical objects but below the Forms. They belong to the
world of Being, but mathematicians still rely on hypotheses and diagrams.

\subsection*{Teaching Point}

Mathematics is elevated above ordinary sensory knowledge, but it is not yet the highest
form of knowledge. It prepares the soul for dialectic and the grasp of the Forms.

\subsection*{Whiteboard}

\[
\begin{array}{c}
\text{GOOD} \\
\hline
\text{FORMS} \\
\hline
\text{MATHEMATICAL OBJECTS} \\
\hline
\text{PHYSICAL OBJECTS} \\
\hline
\text{IMAGES}
\end{array}
\]

\[
\text{mathematics} = \text{middle step upward}
\]

\section*{6. Plato's Criticism of Mathematical Language}

Plato criticizes ordinary geometrical language because it sounds dynamic and constructive.

Geometers say:

\[
\text{draw a line}
\qquad
\text{construct a triangle}
\qquad
\text{move a figure}
\]

But if the true objects of geometry are eternal and unchanging, then they are not literally
drawn, constructed, or moved.

\subsection*{Teaching Point}

Plato is an example of philosophy-first. His metaphysics of mathematical objects leads him
to criticize ordinary mathematical language.

\subsection*{Whiteboard}

\[
\text{Mathematical objects are eternal and unchanging.}
\]

\[
\therefore
\quad
\text{``draw a line'' is philosophically misleading.}
\]

\section*{7. Plato on Arithmetic and Pure Units}

Plato applies the same contrast to arithmetic.

Ordinary arithmetic counts physical things:

\[
\text{two cows}
\qquad
\text{two armies}
\qquad
\text{two shoes}
\]

But these units are unequal. Two shoes differ in shape. Two cows differ in size. Two armies
differ enormously.

Philosophical arithmetic concerns pure units. Each unit is precisely equal to every other
unit.

\[
\begin{array}{c|c}
\textbf{Ordinary arithmetic} & \textbf{Philosophical arithmetic} \\
\hline
\text{physical units} & \text{pure units} \\
\text{unequal} & \text{equal} \\
\text{practical counting} & \text{knowledge of number itself}
\end{array}
\]

\subsection*{Teaching Point}

For Plato, arithmetic, like geometry, is strictly about intelligible objects rather than physical
ones.

\subsection*{Whiteboard}

\[
\text{ordinary two}
=
\text{two unequal physical things}
\]

\[
\text{philosophical two}
=
\text{two pure equal units}
\]

\section*{8. Mathematics on Plato}

The chapter then reverses direction. It asks not only what Plato thought about
mathematics, but how mathematics shaped Plato's philosophy.

Socrates and Plato differ sharply.

\[
\begin{array}{c|c}
\textbf{Socrates} & \textbf{Plato} \\
\hline
\text{public questioning} & \text{elite training} \\
\text{ethics and politics} & \text{mathematics and metaphysics} \\
\text{fallible dialectic} & \text{rigorous demonstration} \\
\text{philosophy for everyone} & \text{philosophy for Guardians}
\end{array}
\]

Socrates uses questioning to expose confusion. Plato, influenced by mathematics, wants
certainty, proof, and disciplined ascent toward first principles.

\subsection*{Teaching Point}

Greek mathematics helped transform philosophy into rationalism: the search for certain
knowledge through reason.

\subsection*{Whiteboard}

\[
\text{Socratic questioning}
\quad \longrightarrow \quad
\text{fallible testing}
\]

\[
\text{Mathematical proof}
\quad \longrightarrow \quad
\text{certainty}
\]

\[
\text{Plato: philosophy should imitate mathematics.}
\]

\section*{9. Mathematics and the Guardians}

In the \emph{Republic}, the future Guardians study mathematics for ten years, from roughly
age 20 to age 30.

This training is not mainly practical. It is spiritual and intellectual. Mathematics turns the
soul away from the world of change and toward reality.

After mathematics, the Guardians study dialectic. Dialectic aims at unhypothetical first
principles and ultimately at the Form of the Good.

\subsection*{Teaching Point}

For Plato, mathematics is not the final goal of philosophy. It is preparation for dialectic.

\subsection*{Whiteboard}

\[
\text{Mathematics}
\quad \longrightarrow \quad
\text{Dialectic}
\quad \longrightarrow \quad
\text{The Good}
\]

\[
\text{hypotheses}
\quad \longrightarrow \quad
\text{first principles}
\]

\section*{10. Aristotle's Rejection of Separate Forms}

Aristotle rejects Plato's separate world of Forms.

For Plato, Forms exist apart from physical things. For Aristotle, Forms exist in things.

\[
\begin{array}{c|c}
\textbf{Plato} & \textbf{Aristotle} \\
\hline
\text{Forms separate from things} & \text{Forms in things} \\
\text{Being apart from Becoming} & \text{form embedded in nature} \\
\text{mathematical objects separate} & \text{mathematics studies aspects of things}
\end{array}
\]

Beauty is not a separate entity apart from beautiful things. Beauty is what beautiful things
have in common.

Likewise, mathematical objects do not need to belong to a separate realm.

\subsection*{Teaching Point}

Aristotle brings Forms down into the world. This changes the philosophy of mathematics:
mathematics must now be understood in relation to perceptible objects.

\subsection*{Whiteboard}

\[
\textbf{Plato:}
\quad
\text{Form apart from things}
\]

\[
\textbf{Aristotle:}
\quad
\text{form in things}
\]

\section*{11. Aristotle on Mathematical Abstraction}

Aristotle says that the mathematician studies physical objects, but not as physical.

A geometer may consider a brass sphere, but not as brass. She considers it as spherical.

\[
\text{brass sphere}
\quad \longrightarrow
\quad
\text{ignore brass}
\quad \longrightarrow
\quad
\text{sphere}
\]

An arithmetician may consider a flock of sheep, but not as sheep. She considers them as
units.

\[
\text{five sheep}
\quad \longrightarrow
\quad
\text{ignore sheepness}
\quad \longrightarrow
\quad
5
\]

The separation is not metaphysical. It is separation in thought.

\subsection*{Teaching Point}

For Aristotle, mathematical abstraction does not require a separate mathematical realm.
The mathematician focuses on certain features of physical things and ignores others.

\subsection*{Whiteboard}

\[
\text{Mathematics studies physical things}
\]

\[
\text{but not qua physical.}
\]

\[
\text{separation in thought}
\neq
\text{separation in reality}
\]

\section*{12. Two Interpretations of Aristotle}

The chapter discusses two interpretations of Aristotle's account.

\[
\begin{array}{c|c}
\textbf{Abstractionist interpretation} & \textbf{Fictionalist interpretation} \\
\hline
\text{mathematical objects are abstracted} & \text{mathematical objects are useful fictions} \\
\text{objects exist as forms of things} & \text{strictly, no separate objects} \\
\text{realism in ontology} & \text{anti-realism in ontology} \\
\text{realism in truth-value} & \text{realism in truth-value}
\end{array}
\]

On the abstractionist view, mathematical objects exist as abstractions from physical objects.

On the fictionalist view, we do not literally create or discover separate mathematical
objects. We pretend there are such objects because doing so is useful and harmless.

\subsection*{Teaching Point}

Both interpretations try to preserve the objectivity of mathematics without Plato's
separate world of Being.

\subsection*{Whiteboard}

\[
\text{Aristotle:}
\quad
\text{objective mathematics without separate Forms?}
\]

\[
\begin{array}{c|c}
\text{Abstraction} & \text{Useful fiction}
\end{array}
\]

\section*{13. Frege's Criticism of Abstraction}

The chapter uses Frege to raise a problem for abstraction.

Suppose we try to obtain the number of a group of objects by abstracting away their
differences. If abstraction removes all differences, then we cannot count them as many.

If the objects become completely identical, then we have only one object.

\[
\text{No difference}
\quad \Rightarrow \quad
\text{no counting beyond one}
\]

But if the differences remain, then abstraction has not produced pure identical units.

\subsection*{Teaching Point}

Abstraction is difficult to explain. If we abstract too much, we lose multiplicity. If we
abstract too little, we do not get pure mathematical objects.

\subsection*{Whiteboard}

\[
\text{Too much abstraction}
\quad \Rightarrow \quad
\text{one indistinguishable thing}
\]

\[
\text{Too little abstraction}
\quad \Rightarrow \quad
\text{ordinary physical differences remain}
\]

\section*{14. Aristotle's Advantage: Applicability}

Aristotle has an advantage over Plato: he explains why mathematics applies to the physical
world.

If mathematics studies aspects of physical things, then it is no mystery that mathematics
applies to physical things.

\[
\text{Mathematics studies physical things under abstraction.}
\]

\[
\therefore
\quad
\text{Mathematics applies to physical things.}
\]

Unlike Plato, Aristotle does not need to explain how a separate mathematical realm is
connected to the physical realm.

\subsection*{Teaching Point}

Aristotle makes applicability straightforward by refusing to separate mathematics from the
physical world.

\subsection*{Whiteboard}

\[
\textbf{Aristotle's advantage:}
\]

\[
\text{No gap between mathematics and nature.}
\]

\[
\text{Application is built in.}
\]

\section*{15. Aristotle's Problem: Exactness}

Aristotle's problem is the reverse of Plato's.

Physical objects are imperfect. A brass sphere is not perfectly spherical. A drawn line has
thickness. A physical tangent overlaps a physical circle over a region.

So if mathematics is about physical things, or abstractions from physical things, how can
mathematics be exact?

\[
\text{imperfect physical object}
\quad \longrightarrow ? \quad
\text{perfect mathematics}
\]

This is the central difficulty for Aristotle.

\subsection*{Teaching Point}

Plato explains exactness but struggles with application. Aristotle explains application but
struggles with exactness.

\subsection*{Whiteboard}

\[
\begin{array}{c|c}
\textbf{Plato} & \textbf{Aristotle} \\
\hline
\text{exactness explained} & \text{application explained} \\
\text{application mysterious} & \text{exactness mysterious}
\end{array}
\]

\section*{16. Limits of Aristotle's Approach}

Aristotle's approach works best for mathematics close to ordinary physical experience:
geometry, counting, ratios, and perhaps parts of analysis.

But modern mathematics extends far beyond direct abstraction from physical objects.

Examples:

\[
\text{complex analysis}
\]

\[
\text{functional analysis}
\]

\[
\text{point-set topology}
\]

\[
\text{axiomatic set theory}
\]

A modern Aristotelian must explain how these areas relate to the physical world, if they
do at all.

\subsection*{Teaching Point}

Aristotle anticipates empiricist approaches to mathematics, but his view faces pressure
from highly abstract modern mathematics.

\subsection*{Whiteboard}

\[
\text{Counting sheep}
\quad \text{and}
\quad
\text{brass spheres}
\]

\[
\neq
\]

\[
\text{set theory, topology, functional analysis}
\]

\section*{Conclusion: Plato and Aristotle as Opposing Models}

The chapter presents Plato and Aristotle as the first major contrast in philosophy of
mathematics.

\[
\begin{array}{c|c}
\textbf{Plato} & \textbf{Aristotle} \\
\hline
\text{mathematics belongs to Being} & \text{mathematics abstracts from Becoming} \\
\text{objects are separate and ideal} & \text{forms are in things} \\
\text{knowledge by reason/recollection} & \text{knowledge through abstraction} \\
\text{exactness is secure} & \text{applicability is secure} \\
\text{application is difficult} & \text{exactness is difficult}
\end{array}
\]

\section*{Final Framing}

To teach this chapter well, present it as a foundational opposition.

Plato and Aristotle are not merely historical figures. They introduce two permanent
tendencies in the philosophy of mathematics.

\[
\boxed{
\text{Platonist tendency: protect exactness by separating mathematics from nature.}
}
\]

\[
\boxed{
\text{Aristotelian tendency: protect applicability by embedding mathematics in nature.}
}
\]

Later philosophies of mathematics repeatedly return to this tension.

\section*{Discussion Questions}

\begin{enumerate}[leftmargin=2em]
    \item Why did Greek mathematicians care so much about exactness rather than approximation?
    \item Does the tangent theorem concern the drawn circle or the ideal circle?
    \item What role do diagrams play in geometrical proof?
    \item Is Plato right that mathematical knowledge is independent of the senses?
    \item Why does Plato think mathematics prepares the soul for philosophy?
    \item Is philosophical arithmetic different from ordinary counting?
    \item Can Aristotle explain exact mathematical truth without a separate mathematical realm?
    \item Is abstraction a plausible source of mathematical objects?
    \item Does Aristotle explain the applicability of mathematics better than Plato?
    \item Which is more important for a philosophy of mathematics: exactness or applicability?
\end{enumerate}

\section*{Key Terms}

\begin{description}[leftmargin=2em]
    \item[World of Becoming] Plato's name for the changing, physical, sensory world.

    \item[World of Being] Plato's name for the eternal, intelligible realm of Forms and mathematical objects.

    \item[Form] A perfect, intelligible standard such as Beauty itself, Justice itself, or Circle itself.

    \item[Rationalism] The view that reason is a fundamental source of knowledge.

    \item[Recollection] Plato's doctrine that learning is remembering what the soul already knew.

    \item[Divided line] Plato's image in the \emph{Republic} representing levels of reality and cognition.

    \item[Mathematical object] An object studied by mathematics, such as a number, point, line, circle, or sphere.

    \item[Pure unit] A unit in philosophical arithmetic, precisely equal to every other unit.

    \item[Dialectic] For Plato, the highest philosophical method, aimed at first principles and the Forms.

    \item[Abstraction] The process of focusing on some features of things while ignoring others.

    \item[Qua] Latin for ``as''; to study something qua mathematical is to study it under its mathematical aspect.

    \item[Fictionalism] The view that mathematical objects are useful fictions rather than literally existing objects.

    \item[Applicability problem] The problem of explaining why mathematics applies to the physical world.
\end{description}

\end{document}